\chapter{Introduction to AI coding agents}

\section{Engineering principles}

A coding assistant suggests a fragment; a coding agent can inspect a workspace, propose edits, invoke tools and respond to results. Traditional engineering retains human control over each action. Agentic engineering delegates a bounded objective while retaining accountability for the resulting change. Building with agents means applying existing execution platforms to software delivery; building agents means creating a new agent application; building harnesses means engineering the control plane that surrounds inference. This book uses the first interpretation and extends existing harnesses only where organizational requirements demand it.

\section{Implementation approach}

Select tasks by ambiguity, blast radius and testability. Start with a small deterministic defect, a documentation update or an isolated refactor. Avoid delegating production administration, credential rotation, irreversible data changes or open-ended architectural redesign without approvals. A task owner specifies the objective, prohibited actions, permitted files and observable completion evidence. Use the baseline repository commit as the comparison point.

\section{Verification and limitations}

The engineer remains accountable for correctness, risk and review. Measure success by a reviewed patch with independent evidence, not by generated lines of code or a confident completion message. Reject changes that cannot be traced to a requirement.

\section{Chapter conclusion}

Apply the stated verification criteria before expanding the scope of this workflow. Record deviations, evidence and unresolved risks in the task record.
